\documentclass[11pt,a4paper,fontset=fandol]{ctexart}

\usepackage[a4paper,top=2.25cm,bottom=2.45cm,left=2.25cm,right=2.25cm,headheight=18pt,headsep=0.55cm,footskip=24pt]{geometry}
\usepackage{amsmath,amssymb,amsthm}
\usepackage{fancyhdr}
\usepackage{lastpage}
\usepackage{enumitem}
\usepackage{titlesec}
\usepackage{xcolor}
\usepackage{hyperref}
\usepackage{bookmark}
\usepackage[most]{tcolorbox}
\usepackage{setspace}
\usepackage{array}

\hypersetup{
  colorlinks=true,
  linkcolor=blue!50!black,
  urlcolor=blue!50!black
}

\setstretch{1.13}
\setlength{\parindent}{2em}
\setlength{\parskip}{0.25em}
\allowdisplaybreaks
\setlist[enumerate]{itemsep=0.45em, topsep=0.35em, leftmargin=2.4em}
\setlist[itemize]{itemsep=0.25em, topsep=0.25em, leftmargin=1.9em}

\definecolor{mainblue}{RGB}{36,83,140}
\definecolor{lightblue}{RGB}{241,246,252}
\definecolor{lightgray}{RGB}{246,246,246}
\definecolor{rulegray}{RGB}{110,118,128}

\titleformat{\section}
  {\Large\bfseries\color{mainblue}}
  {\thesection.}
  {0.45em}
  {}
  [\vspace{0.2em}\color{rulegray}\titlerule]
\titlespacing*{\section}{0pt}{1.1em}{0.7em}

\pagestyle{fancy}
\fancyhf{}
\fancyhead[L]{\small 函数图像综合变换周测 A 卷}
\fancyhead[C]{\small 代数专题：基础与标准变式}
\fancyhead[R]{\small 刘文博}
\fancyfoot[C]{\small 第 \thepage 页 / 共 \pageref{LastPage} 页}
\renewcommand{\headrulewidth}{0.55pt}
\renewcommand{\footrulewidth}{0.35pt}

\newtcolorbox{infobox}[1]{
  breakable,
  colback=lightblue,
  colframe=mainblue,
  boxrule=0.7pt,
  arc=1mm,
  title=\textbf{#1},
  fonttitle=\bfseries
}

\newenvironment{solution}{\par\noindent\textbf{参考解答：}}{\hfill$\square$\par}
\newcommand{\answerspace}{\vspace{2.3cm}}
\newcommand{\biganswerspace}{\vspace{3.1cm}}

\begin{document}

\thispagestyle{empty}
\begin{center}
{\fontsize{22}{28}\selectfont\bfseries 函数图像综合变换周测 A 卷\par}
\vspace{0.4cm}
{\large 适合初中阶段函数专题基础检测\par}
\end{center}

\begin{infobox}{考试说明}
\begin{itemize}
  \item 测试时间：60--70 分钟
  \item 满分：100 分
  \item A 卷侧重标准综合变换、反向判断和关键点跟踪，先保证分层改式子稳定。
\end{itemize}
\end{infobox}

\section{试题}

\begin{enumerate}
  \item （10 分）把函数
  \[
  y=x^2
  \]
  先向左平移 $2$ 个单位，再关于 \(x\) 轴对称，求新方程。
  \answerspace

  \item （12 分）把函数
  \[
  y=2x+1
  \]
  先关于原点对称，再向上平移 $4$ 个单位，求化简后的新方程。
  \answerspace

  \item （12 分）把函数
  \[
  y=|x|
  \]
  先关于 \(y\) 轴对称，再向右平移 $3$ 个单位，求新方程。
  \answerspace

  \item （12 分）已知函数
  \[
  y=-f(x-2)+5,
  \]
  说出它是由
  \[
  y=f(x)
  \]
  经过哪些变换得到的。
  \answerspace

  \item （12 分）比较下列两个结果是否一定相同，并说明理由：
  \begin{enumerate}
    \item 先向右平移 $2$ 个单位，再关于 \(x\) 轴对称；
    \item 先关于 \(x\) 轴对称，再向右平移 $2$ 个单位。
  \end{enumerate}
  \answerspace

  \item （14 分）已知点 \((1,3)\) 在函数
  \[
  y=f(x)
  \]
  的图像上。若新图像由原图像先向左平移 $4$ 个单位，再关于 \(x\) 轴对称得到，求对应点的新坐标。
  \biganswerspace

  \item （14 分）已知抛物线
  \[
  y=x^2
  \]
  经过综合变换后，顶点变成 \((3,2)\)，开口方向向下。求一个对应的新方程。
  \biganswerspace

  \item （14 分）把函数
  \[
  y=|x-1|+2
  \]
  先关于 \(y\) 轴对称，再向下平移 $3$ 个单位，求新方程，并写出新图像的顶点坐标。
  \biganswerspace
\end{enumerate}

\newpage
\section{详细解答}

\begin{enumerate}
  \item \begin{solution}
  先向左平移 $2$ 个单位，得到
  \[
  y=(x+2)^2.
  \]
  再关于 \(x\) 轴对称，得到
  \[
  \boxed{y=-(x+2)^2}.
  \]
  \end{solution}

  \item \begin{solution}
  先关于原点对称：
  \[
  y=-(2(-x)+1)=2x-1.
  \]
  再向上平移 $4$ 个单位：
  \[
  y=2x-1+4=2x+3.
  \]
  所以新方程为
  \[
  \boxed{y=2x+3}.
  \]
  \end{solution}

  \item \begin{solution}
  函数 \(y=|x|\) 本身关于 \(y\) 轴对称，所以第一步以后仍为
  \[
  y=|x|.
  \]
  再向右平移 $3$ 个单位，得到
  \[
  \boxed{y=|x-3|}.
  \]
  \end{solution}

  \item \begin{solution}
  在方程
  \[
  y=-f(x-2)+5
  \]
  中：
  \begin{itemize}
    \item \(x-2\) 表示向右平移 $2$ 个单位；
    \item 最前面的负号表示关于 \(x\) 轴对称；
    \item 最后的 \(+5\) 表示向上平移 $5$ 个单位。
  \end{itemize}
  因此它是由原图像先向右平移 $2$ 个单位，再关于 \(x\) 轴对称，最后向上平移 $5$ 个单位得到的。
  \end{solution}

  \item \begin{solution}
  一定相同。

  过程 1：
  \[
  y=f(x) \rightarrow y=f(x-2) \rightarrow y=-f(x-2).
  \]
  过程 2：
  \[
  y=f(x) \rightarrow y=-f(x) \rightarrow y=-f(x-2).
  \]

  两个结果相同，因为一个改横坐标，一个改函数值，彼此不冲突。
  \end{solution}

  \item \begin{solution}
  点 \((1,3)\) 先向左平移 $4$ 个单位，变成 \((-3,3)\)。

  再关于 \(x\) 轴对称，变成 \((-3,-3)\)。

  所以对应点的新坐标是
  \[
  \boxed{(-3,-3)}.
  \]
  \end{solution}

  \item \begin{solution}
  顶点从 \((0,0)\) 变成 \((3,2)\)，说明图像向右平移 $3$ 个单位，再向上平移 $2$ 个单位；开口向下说明还要关于 \(x\) 轴对称。

  所以一个对应的新方程为
  \[
  \boxed{y=-(x-3)^2+2}.
  \]
  \end{solution}

  \item \begin{solution}
  先关于 \(y\) 轴对称：
  \[
  y=|x+1|+2.
  \]
  再向下平移 $3$ 个单位：
  \[
  y=|x+1|-1.
  \]
  所以新方程为
  \[
  \boxed{y=|x+1|-1}.
  \]

  原顶点 \((1,2)\) 关于 \(y\) 轴对称后变成 \((-1,2)\)，再向下平移 $3$ 个单位，得到
  \[
  \boxed{(-1,-1)}.
  \]
  \end{solution}
\end{enumerate}

\end{document}
